\section{Thirteen integers: CRT elimination and a coprime spectator guard}
\label{sec:jones-seeded}

The polynomial specified below has exactly thirteen unrestricted integer
unknowns, exact total degree $28\,847\,344$, and a division-free circuit of
$754$ live binary operations ($516$ multiplications, $215$ additions and
$23$ subtractions). Its integer zero existence is equivalent to a finite
syntactic ZFC refutation. Thus, with the stated inherited compiler theorem,
ZFC proves that its lack of integer zeros is equivalent to
$\operatorname{Con}(\mathrm{ZFC})$. This is an ordinary mathematical proof,
not a new Lean certificate, a proof of consistency, or a global minimality
claim. Unlike the ten-integer companion, the arithmetic finalizer here uses
only the established Jones ratio interface, not Sun's 2026 bridge.

\subsection{A coherent fixed source and an explicit polynomial}
Let $p_1,\ldots,p_{22}$ be the original package-31 quadratic residuals after
substituting $x_i=n_i+1$ in all $31$ positive source coordinates. The exact
source is \texttt{upstream/parent\_quartic31.json}; the complete fixed port
mapping is $(\beta,K_0,M_C,M_F)=(\beta,K,M_C,M_{F,\mathrm{source}})$ from
\path{FIXED_BINDING_OLD_SOURCE.json}. All four compressed little-endian payloads are
included and hash-authenticated. This uses the original transport coordinate,
not package 45's extra transport shift. The imported source theorem says
\[
 \exists n\in\mathbb N^{31}\ \bigwedge_{j=1}^{22}p_j(n)=0
 \quad\Longleftrightarrow\quad
 \text{there is a finite ZFC refutation}.
\]
The fifth residual equals $1$ at $n=0$, so any common zero is nonzero.

Use the selective weights $w_i$, row order, and targets $q_j$ in
\texttt{CODE.json}, identical to the companion's certified tuple.
Their maxima are $W=183$ and $Q_*=5190$; put $K_*=Q_*+2W=5556$ and
$E_*=K_*+W+2=5741$. For
$p_j(n)=\sum_\alpha p_{j,\alpha}n^\alpha$, set
\[
 c_{j,\alpha}=(2-|\alpha|)!\prod_i\alpha_i!\,p_{j,\alpha},\qquad
 D_c(B)=\sum_{j,\alpha}c_{j,\alpha}B^{q_j-\alpha\cdot w}.
\]
The finite isolation certificate checks all $71\,808$ possible interactions
of the $136$ used source terms with the $528$ padded quadratic monomials.
It proves
\[
 [t^{q_j}]\Bigl(1+\sum_i n_it^{w_i}\Bigr)^2D_c(t)=2p_j(n).
\]
It also proves that the displayed exponents in $D_c$ are distinct and its
leading coefficient is $2$. All $|c_{j,\alpha}|<A_*=2^{2134927}$;
the verifier derives this bound from the actual fixed-port bit lengths.

The unknowns are
\[
 (\epsilon,g,h,s,w,\phi,i,j,z_1,z_2,z_3,r,v)\in\mathbb Z^{13}.
\]
Every subsequent symbol is a polynomial register, not an additional unknown.
Fix
\[
 \kappa=2134940,\quad a=2^\kappa=8192A_*,\quad
 \eta=2^{16012040}=2^{15\kappa/2-10},
\]
and define
\begin{align*}
 b&=2\epsilon^2,& B&=ab^2,\\
 \ell&=\sum_iB^{w_i},& T_*&=\sum_jB^{q_j},\\
 N_0&=B^{W+1},& L&=2bB^{E_*},\\
 N&=\eta B^{2966},& c&=bB^W,\\
 M_c&=N_0-1-(b-1)\ell,&
 J_c&=2(1+g)^2D_c(B)+T_*,\\
 S_c&=g+N_0J_c,& T_c&=M_c+N_0(B-2)T_*,\\
 R&=L^2+(S_c+T_c)L+(L-1-S_c),& G_c&=32c.
\end{align*}
All powers have fixed nonnegative integer exponents. In particular, $N$ is
an ordinary polynomial, not a variable-exponent abbreviation.

Use the Jones registers with the exact CRT reparameterization of $H$ below, and $\omega=w^2$:
\begin{align*}
 U&=N^2\omega,&Y&=N^2s,\\
 M&=RY,&P&=2M^2U,\\
 K&=R+1+h(P-1),&A&=M(U+1),\\
 B_1&=2R+1,&W_1&=b\omega,\\
 C&=B_1+W_1+\phi,&J&=4A-5,&
 Q&=3W_1C-2(W_1^2-1),\\
 D&=(A^2-1)C^2+1,&E&=iJDC^2,&F&=(A^2-1)E^2+1,\\
 G&=A+F(CD+1-A),&H&=C+F(B_1+jC),&I&=(G^2-1)H^2+1,\\
 a_0&=(P^2-1)K^2+1,&a_1&=DFI,\\
 u&=J,&t&=Q.
\end{align*}
Set
\[
 \delta=C-KY,\qquad \Delta=g(c-g)K^2-c^2\delta^2,
 \qquad (d_1,d_2,d_3)=(3s-1,3\phi+2,3\Delta-1).
\]
The new coprime-spectator stack is
\[
 C_0=a_0,\qquad B_j^*=1+4d_jC_{j-1}z_j^2,\qquad
 C_j=C_{j-1}B_j^*\quad(j=1,2),
\]
\[
 B_3^*=1+4d_3C_2a_1z_3^2,\qquad C_*=C_3=C_2B_3^*.
\]
Thus the last guard contains the second square candidate $a_1$ as a
spectator. This is the only change to the Pell sign stack.
Finally define
\begin{align*}
 Z&=\epsilon wi z_1z_2z_3a_0,&\Sigma&=Zu^5,&\rho&=t+ur,\\
 \Theta&=(\rho^3-u^2\rho-u^3)(\rho^2+23u^2),&
 X&=\Sigma v-\Theta,\\
 A_0&=\Sigma^2C_*,&A_1&=\Sigma^2a_1.
\end{align*}
The requested thirteen-integer equation is the following single polynomial:
\begin{equation}\label{eq:jones-seeded}
 \boxed{\mathcal F_{13}=(X^2-A_0-A_1)^2-4A_0A_1=0.}
\end{equation}
The literal division-free DAG is \texttt{generated/signed\_coprime\_crt.json}.

\subsection{Sparse padding and unconditional size estimates}
Assume first only $b\ge2$. Since $B-1>A_*$ and the top coefficient of $D_c$
is $2$, we have $D_c(B)>B^{Q_*}>0$. Also $0<M_c<N_0$.
In fact $S_c>0$ for every integer $g$, not merely for $g\ge1$.
Put $v_g=g+1$ and $A_{\rm sc}=2N_0D_c\ge2$. Then
\[
 S_c=A_{\rm sc}v_g^2+v_g+N_0T_*-1>0,
\]
because $A_{\rm sc}v_g^2+v_g\ge0$ for every integer $v_g$, and $N_0T_*>1$.
Also $T_c>0$. Consequently, unconditionally in $g$,
\[
 R=L^2+S_c(L-1)+T_cL+L-1>L^2.
\]
Moreover
\[
 N^2=2^{-20}a^{15}B^{5932}<B^{5947}<B^{11482}<L^2,
 \qquad b\le c\le N.
\]
Thus $R\ge N\ge8$ and $b\le N,R$ before assuming any carry or square
condition. Negative powers here occur only in this comparison proof, never
in the equation.

If $1\le g<4bB^W$, integrality gives $1+g\le4bB^W$ and $g<N_0$.
The injective coefficients imply
\[
 |D_c(B)|<\frac{A_*B^{Q_*+1}}{B-1},\qquad
 2(1+g)^2D_c(B)<\frac{32A_*b^2}{B-1}B^{K_*+1}.
\]
Because $B=8192A_*b^2$, the last prefactor is less than $1/128$.
Together with $T_*<B^{Q_*+1}/(B-1)$ this gives
$0<J_c<B^{K_*+1}$. Likewise $(B-2)T_*<B^{Q_*+1}\le B^{K_*+1}$.
Consequently
\[
 0<S_c,T_c<B^{E_*},\qquad S_c+T_c<2B^{E_*}<L.
\]
These estimates hold for arbitrary integer $b\ge2$, not just powers of two.

Now suppose $b=2^k$. Then $B=2^\nu$, where $\nu=\kappa+2k$.
The binary mask $M_c$ forces exactly
\[
 g\mathbin{\&}M_c=0\quad\Longleftrightarrow\quad
 g=\sum_i n_iB^{w_i},\qquad 0\le n_i<b.
\]
Here $\&$ is used only in the proof. For such a code, write
$h_e=[t^e](1+\sum n_it^{w_i})^2D_c(t)$. Then
\[
 |h_e|\le A_*(1+\textstyle\sum_i n_i)^2
 <1024A_*b^2=B/8<B/4.
\]
The unnormalized digits of $J_c$ are $2h_e+\varepsilon_e$, where
$\varepsilon_e=1$ at a target and $0$ elsewhere. Starting with incoming
carry $c_0=0$, normalized carries obey
\[
 c_{e+1}=\left\lfloor\frac{2h_e+\varepsilon_e+c_e}{B}\right\rfloor,
 \qquad c_e\in\{-1,0\}.
\]
At a target, its normalized digit avoids $B-2$ exactly when it is $0$ or $1$.
If $h_e=0$, the digit is $1+c_e\in\{0,1\}$. If $h_e>0$, it is at least $2$;
if $h_e<0$, a borrow normalizes it to a digit at least $2$.
Thus the target test is equivalent to $h_e=0$, independently of incoming
borrows. There is no need to fill every position with a dense $B/2$ offset.
It follows, by splitting at $N_0$, that
\[
 S_c\mathbin{\&}T_c=0\quad\Longleftrightarrow\quad
 g\text{ codes a common zero of all }p_j.
\]
Conversely every common zero has, for arbitrarily large $b=2^{1+2r}$, such a
code with $1\le g<bB^W=c$.

\subsection{The parity-tuned two-block threshold}
Write $L=2^\lambda$ after $b=2^k$ has been established. The three disjoint
blocks in $R$ give
\[
 s_2(R)=1+s_2(S_c+T_c)+\lambda-s_2(S_c)
       =\lambda+s_2(T_c)+1-\tau_2(S_c,T_c),
\]
where $s_2$ is binary digit sum and
$\tau_2(S,T)=s_2(S)+s_2(T)-s_2(S+T)\ge0$ counts binary carries.
The strict inequality $S_c+T_c<L$ is essential.
The $31$ allowed low positions and the $22$ target positions give
\[
 s_2(T_c)=(W+23)\nu-31k-22,\qquad
 \lambda=1+k+E_*\nu.
\]
Therefore
\begin{align*}
 \lambda+s_2(T_c)+1
 &=5947\kappa+11864k-20\\
 &=2\bigl(15\kappa/2-10+2966(\kappa+2k)\bigr)
 =\log_2N^2.
\end{align*}
Legendre's formula yields
$v_2\binom{2R}{R}=s_2(R)$, hence
\[
 N^2\mid\binom{2R}{R}
 \quad\Longleftrightarrow\quad S_c\mathbin{\&}T_c=0.
\]
The small factor $2b$ in $L$ calibrates the threshold to a polynomial square
without an extra variable, an extra target mask, or a cubic radix.

\subsection{An exact polynomial CRT reparameterization}
The inherited ratio interface uses $H_{\rm old}=B_1+j_{\rm old}C$ and
$JF\mid QF+J(H_{\rm old}-C)$. Since $E=iJDC^2$,
\[
 F\equiv1\pmod J,\qquad F\equiv1\pmod C,
\]
and hence $\gcd(J,F)=1$. The divisibility is equivalent to
$J\mid Q$ and $F\mid H_{\rm old}-C$. Our changed register
\[
 H=C+F(B_1+jC)
\]
automatically satisfies the latter. It also has the original form, with
\[
 j_{\rm old}=1+(A^2-1)i^2J^2D^2C^3B_1+Fj,
\]
which is a polynomial integer, since $F-1=(A^2-1)i^2J^2D^2C^4$.
This identity holds for arbitrary signed inputs, including $C=0$.

Conversely, at the canonical positive ratio witnesses, $0<B_1<C$ and
$j_{\rm old}\ge1$. Therefore $z=(H_{\rm old}-C)/F>0$ and
$z\equiv B_1\pmod C$, so $z=B_1+jC$ with $j\ge0$. Thus every
canonical witness is represented. No extra unknown or cofinal restriction is
introduced. The merger now has $u=J$ rather than $JF$. This CRT step is
inherited from the separately identified sparse-Pell packet, not claimed as
an original theorem here.

\subsection{The interval guard and the Jones interface}
If $\Delta>0$ and $b\ge2$, then $K\ne0$ and
$0<g<c$. The identity $4g(c-g)=c^2-(c-2g)^2\le c^2$ implies
$4\delta^2<K^2$. Along with $s\ge1$, $\phi\ge0$, and $w\ne0$, this gives
$C,Y>0$ and $K>0$. The usual square-gap argument forces $h\ge1$:
if $h\le0$, then $1\le K\le R+1<P$; $K=1$ would force $P-1\mid R$,
which is impossible, while for $2\le K<P$,
\[
 (PK-1)^2<(P^2-1)K^2+1<(PK)^2.
\]
Thus $a_0$ cannot be a square in that branch.

For precision, instantiate the free ceiling in
\texttt{Jones1982.Nine.ratio\_sound} by
$G_{\mathrm{aux}}=8gK^2+1$, rather than assuming its old margin condition.
The positive integer $K^2-4\delta^2$ then gives
\[
 (K^2-4\delta^2)G_{\mathrm{aux}}-8gK^2>0.
\]
This ceiling is a proof-side abbreviation; the ratio registers do not
otherwise depend on it.

Replace $i$ by $|i|$, which changes no condition. To handle signed $j$,
$u\mid t$ now means $J\mid Q$, while $F\mid H-C$ holds identically. Also
$F\equiv1\pmod D$, $G\equiv1\pmod D$, and $G\equiv A\pmod F$ show
that the positive integers $D,F,I$ are pairwise coprime.
Hence $a_1=DFI$ square makes each a square. Put $d=G^2-1>0$ and
$x=\sqrt I$. The norm-one unit $x+H\sqrt d$ is positive even if $H<0$.
A positive power of $G+\sqrt d$ is the identity modulo $CF$; multiplying by
sufficiently large such powers gives a Pell solution $H'>B_1$ with
$H'\equiv H\pmod{CF}$. Then $j'=(H'-B_1)/C>0$, all square conditions
remain true, and $J\mid Q$ remains true. The literal positive Jones interface
therefore applies and gives the power-of-two base and central-binomial
condition required above.

For completeness, take any common source zero, choose
$b=2^{1+2r}$ large enough, and set $\epsilon=2^r$. The direct source code
has $32g<G_c$, so the literal ratio completeness theorem supplies its
canonical witnesses. In \texttt{Ratio.lean}, \texttt{exists\_B25core} sets
\[
 \omega=2^{2R+1-k},\qquad U=(N^2/b)2^{2R+1}.
\]
Since $k=1+2r$ is odd and $k<2R+1$, this $\omega$ is a positive square.
No claim that the canonical $s$ is a square is needed.

The same canonical construction has
$\rho_0=(U+1)^{2R}/U^R=Y+f$ with
$0\le f<2^{2R}/U=b/(2N^2)$ and $Y\ge1$.
The source's \texttt{ratio\_core} gives
\[
 |C/K-\rho_0|\le\frac{2(Y+f)}{Y(U+1)}<\frac4U,
 \qquad |C/K-Y|<\frac b{N^2}.
\]
The last inequality uses $R\ge8$. Since $c,b\le N$ and
$g(c-g)\ge1$, it follows that
\[
 \frac{c^2(C-KY)^2}{K^2}<\frac{c^2b^2}{N^4}\le1\le g(c-g),
\]
so the very same canonical witnesses, translated by the exact CRT formula, satisfy $\Delta>0$.
This is a cofinal square-base proof, not an unsupported restriction of an
existing slack variable to squares.

\subsection{Raw two-root recovery, including the zero spectator branch}
For every integer assignment, $u=J=4A-5\ne0$. The polynomial
\[
 \Pi(x)=(x^3-x-1)(x^2+23)
\]
has no rational root and has a root modulo every positive integer.
For an odd prime other than $23$, either the cubic has a simple root or it
is irreducible, whose Frobenius three-cycle makes its nonzero discriminant
$-23$ a square. The quadratic then has a simple root. At $23$ the cubic has
the simple root $3$; at $2$, $-23\equiv1\pmod8$ admits compatible square
roots. Hensel lifting and CRT finish the assertion. Thus
$\Theta=u^5\Pi(t/u+r)\ne0$ for every assignment, and
\[
 \exists r\quad Zu^5\mid\Theta
 \quad\Longleftrightarrow\quad Z\ne0\ \hbox{and}\ u\mid t.
\]
Indeed a root modulo $|Z|$ proves the reverse implication once $u\mid t$.
If $v_p(u)=\alpha>\beta=v_p(t)$, the two factors of $\Theta$ have
valuations $3\beta$ and $2\beta$, contrary to $u^5\mid\Theta$.
If $Z=0$, divisibility would require $\Theta=0$, already impossible.

At a zero of the equation, $\Sigma=0$ is impossible because then the
polynomial is $\Theta^4\ne0$. In particular
$\epsilon,w,i,z_1,z_2,z_3,a_0$ are nonzero. Every $d_j\equiv2\pmod3$
is nonzero, and every $B_j^*\equiv1\pmod4$ is nonzero. Thus $C_2,C_*$
are nonzero.

\paragraph{No coincident radicands.}
One has $C_*\ne a_1$ without any positivity assumption. If $a_1=0$ this
follows from $C_*\ne0$. Otherwise $B_3^*\ne1$, and it cannot be $-1$
because it is $1$ modulo $4$. But
\[
 \gcd(B_3^*,a_1)=1,
\]
while $C_*=a_1$ would make $B_3^*$ divide $a_1$. This would force
$B_3^*=\pm1$, a contradiction.

Put $q=X/\Sigma$. The raw norm equation implies
\[
 (q^2+C_*-a_1)^2=4C_*q^2.
\]
If $q=0$, the norm would give $C_*=a_1$, just excluded. Therefore
\[
 r_0=\frac{q^2+C_*-a_1}{2q},\qquad r_1=q-r_0
\]
are rational, with $r_0^2=C_*$ and $r_1^2=a_1$. Rational square roots of
integers are integers. Hence $C_*,a_1$ are integer squares,
$q\in\mathbb Z$, and $\Sigma\mid\Theta$. In particular $J\mid Q$.
No complex-root separation assumption is used.

\paragraph{Recover the first two guards before ruling out $a_1=0$.}
The factors $a_0,B_1^*,B_2^*,B_3^*$ are pairwise coprime, because each
new factor is $1$ modulo its entire preceding product. Their nonzero product
$C_*$ is a square. Thus each absolute factor is a square. Since
$B_j^*\equiv1\pmod4$, every $B_j^*$ is positive: a negative square cannot
be $1$ modulo $4$. The product sign then forces $a_0>0$, also a square.
All prefixes are now positive, so $d_1,d_2>0$, yielding $s\ge1$ and
$\phi\ge0$.

At this point the last factor might still equal $1$ if $a_1=0$; this branch
must not be silently discarded. The unconditional source estimate above
has already proved $R>L^2>0$ for \emph{every} signed $g$. Since
$\epsilon,w\ne0$ and $s\ge1$, we have $b\ge2$, $N\ge8$, $U,Y>0$,
$A\ge2$, and, using $\phi\ge0$, $C>0$. Then
\[
 D\ge A^2>0,\qquad F\ge1,\qquad
 CD+1-A>0,\qquad G\ge A>1,\qquad I\ge1.
\]
Consequently $a_1=DFI>0$. Now $B_3^*\ne1$, and its positivity with
$C_2a_1z_3^2>0$ forces $d_3>0$, hence $\Delta>0$. This closes the
zero-spectator branch without adding $a_1$ to the nonzero multiplier.
The interval and ratio arguments above therefore apply.

\paragraph{Completeness.}
Start from the canonical positive ratio witnesses, translated by the CRT
parameterization. They have $a_0,a_1$ positive squares, $J\mid Q$, and all
three strict sign conditions. For $j=1,2$, $d_jC_{j-1}$ is positive
nonsquare because $d_j\equiv2\pmod3$ and $C_{j-1}$ is a square.
A nontrivial Pell solution $x^2-d_jC_{j-1}y^2=1$ can be chosen with $y$
even, so $z_j=y/2\ne0$ makes $B_j^*=x^2$. For the last guard,
$d_3C_2a_1$ is likewise positive nonsquare, since $C_2a_1$ is a positive
square. The same construction gives $B_3^*$ square and $C_*$ square.
The intersective lemma chooses $r$ with $\Sigma\mid\Theta$. Finally choose
\[
 v=\Theta/\Sigma+\sqrt{C_*}+\sqrt{a_1}\in\mathbb Z.
\]
Then the raw norm vanishes. Both directions are proved. By the inherited
compiler theorem, integer-zero existence is equivalent to a finite ZFC
refutation, and ZFC proves its negation equivalent to
$\operatorname{Con}(\mathrm{ZFC})$.

\subsection{Exact degree, comparison, and verification scope}
The exact relevant register degrees are
\[
\begin{array}{c|rrrrrr}
 &a_0&a_1&C_2&C_*&\Sigma&X\\\hline
 \deg&652210&6592598&2608849&12185330&1119171&1119172
\end{array}
\]
and $\deg A_0=14423672$, $\deg A_1=8830940$. The degree of $A_0$
strictly exceeds those of $X^2$ and $A_1$. The output has leading term
$A_0^2$. Its diagonal leading coefficient is
\[
 72338\,a^{7211252}\eta^{960}\pmod{1000003},
\]
which is nonzero since $a$ and $\eta$ are powers of two. The circuit upper
bound and this specialization together prove exact degree $28\,847\,344$.

At the identical fixed source and selective code, the previous seeded
construction has degree $57\,866\,184$ and $759$ operations. The coprime
spectator alone gives degree $46\,680\,720$ and $758$ operations; combining
it with CRT gives degree $28\,847\,344$ and $754$ operations, a $50.15\%$
degree reduction against the seeded predecessor. All have thirteen integer
unknowns. This is not a twelve-variable construction, a variable-count
record, a global optimum, or a claim that coefficient storage decreased.

The canonical circuit additionally factors the merger cubic as
$\rho(\rho^2-u^2)-u^3$, saving one multiplication without changing the
polynomial. The metric counts every live binary addition, subtraction and
multiplication, including scalar multiplication and fixed exponentiation.
Reuse and fixed coefficient leaves are free. The code independently replays
literal formulas over three finite fields, checks signed and zero branches,
the norm identity, the CRT identity, nonvacuous Pell witness examples, source
positivity, sparse carries, and the inherited fixed-byte hashes. Those are
finite implementation audits, not a replacement for the universal proof or
a new kernel-checked formal theorem.

Run \texttt{python build\_jones.py}, \texttt{python check\_jones.py}, and
\texttt{python -O check\_jones.py}. The complete equation and fixed binding
are in the emitted DAG and \path{FIXED_BINDING_OLD_SOURCE.json}. Source provenance
is retained in \texttt{SOURCE\_PINS.json}; the CRT source packet and its
archive hash are recorded separately in \texttt{CRT\_PROVENANCE.json}.
